\chapter{Brachytherapy}

\section*{Overview}
Brachytherapy is internal radiation therapy in which a radioactive source is placed inside or next to a tumor for a planned period.

\section*{Why It Is Done}
It treats selected cancers, including prostate, cervical, uterine, breast, skin, and other tumors, either alone or with external-beam radiation or surgery.

\section*{How It Is Performed}
Applicators, needles, or small seeds are positioned using imaging and treatment planning. The source may be temporary and removed after treatment or permanently implanted at a low dose rate.

\section*{Risks and Follow-up}
Side effects depend on the treatment site and include bleeding, infection, urinary or bowel symptoms, fatigue, and damage to nearby tissue. Radiation-oncology follow-up monitors tumor response and late effects.

\section*{References}
\begin{enumerate}
  \item National Cancer Institute. Brachytherapy. 2025.
  \item Current Medical Diagnosis and Treatment. McGraw-Hill; 2025.
\end{enumerate}
\clearpage
